% Pattern: proof/amsthm/amsart. A proof whose head gives the name of a
% statement, as [Proof of computational adequacy], proves the statement of
% that name; "Proof of the theorem" proves the last theorem before it; and a
% table between a statement and its proof does not break them.
% seen in the gold pilot (set X)
% want thm-1 statement/theorem 1
% want lem-1 statement/lemma 2
% want proof-1 proof - proves=lem-1
% want proof-2 proof - proves=thm-1
% want thm-2 statement/theorem 3
% want lem-2 statement/lemma 4
% want proof-3 proof - proves=lem-2
% want proof-4 proof - proves=thm-2
% want lem-3 statement/lemma 5
% want proof-5 proof - proves=lem-3
\documentclass{amsart}
\usepackage{longtable}
\newtheorem{thm}{Theorem}
\newtheorem{lem}[thm]{Lemma}
\begin{document}
\begin{thm}[Computational adequacy] Each closed term of ground type halts. \end{thm}

We first prove a lemma.

\begin{lem} Each value halts. \end{lem}
\begin{proof} By the rules. \end{proof}

Adequacy follows from the lemma.
\begin{proof}[Proof of computational adequacy \textnormal{(see the code)}]
Use the logical relation. \end{proof}

\begin{thm} Each tree with $n$ vertices has $n-1$ edges. \end{thm}

We shall prove this by induction.

\begin{lem} Each tree with two vertices or more has a leaf. \end{lem}
\begin{proof} Take a longest path. \end{proof}

\begin{proof}[Proof of the theorem] Remove a leaf. \end{proof}

\begin{lem} The values of the table are right. \end{lem}

\begin{longtable}{ll}
$n$ & $n-1$ \\
\end{longtable}

The table lists the counts.

\begin{proof} Count the edges. \end{proof}
\end{document}
